\section{Certificate schemas}
\label{app:schemas}

The nine prototypes serialised the same mathematics in nine mutually
incompatible ways. This appendix fixes one canonical shape and records the
invariants a checker must enforce. The differences between the original formats
were confined to key names and one numeric encoding: no information is lost in
either direction, and no float rounding is involved anywhere
(\S\ref{app:schema-divergence}).

\subsection{The universal rule}

\begin{invariant}[Exact rationals only]
Every mathematical coefficient is an exact rational, serialised either as a
string such as \code{"125/18"} or as an integer numerator/denominator pair.
No mathematical value is ever a JSON floating-point number. A decimal is not
parsed into a claim of approximate equality, and there is no epsilon, numerical
norm, or solver-status field anywhere in an acceptance condition.
\end{invariant}

The only floats in the entire corpus are timing diagnostics.

\subsection{Sparse polynomials}

A polynomial carries a variable count and a list of
\code{[exponentVector, rationalString]} pairs:

\begin{lstlisting}
{
  "nvars": 2,
  "terms": [
    [[0, 2], "1"],
    [[1, 1], "-2"],
    [[2, 0], "1"]
  ]
}
\end{lstlisting}

Canonical construction rejects malformed exponent vectors, negative exponents,
zero coefficients, arity mismatches, and \emph{duplicate or unsorted monomials}.
Decoding must also reject duplicate JSON object keys: two parsers disagreeing
about which value wins is a genuine differential hazard, and at least one
prototype installs an explicit duplicate-key hook for this reason.

The variable table binding these indices to Lean expressions is a separate
production concern and is deliberately absent from the Python formats, which use
fixed mathematical variables.

\subsection{Cone and ideal certificates}

\begin{lstlisting}
{
  "kind": "nonnegative_cone",
  "terms": [
    { "weight": "1/2",           -- exact, nonnegative
      "square": {...polynomial}, -- this stores q, NOT q^2
      "factors": [0, 2] }        -- indices into the supplied hypotheses
  ],
  "equalities": [
    { "index": 1,                -- index into the supplied equations
      "multiplier": {...} }      -- UNRESTRICTED sign
  ]
}
\end{lstlisting}

Checking, expressed as the loop every implementation runs:

\begin{lstlisting}
acc := zero_polynomial
for term in certificate.terms:
    reject unless term.weight is an exact rational and >= 0
    reject unless term.square has the target's arity
    p := term.weight * term.square^2
    for i in term.factors:
        reject unless i is a valid hypothesis index
        p := p * nonnegative_hypotheses[i]
    acc := acc + p
for (index, multiplier) in certificate.equalities:
    reject unless index is a valid equality index
    acc := acc + multiplier * equalities[index]
accept exactly when acc == target
\end{lstlisting}

Four points that the prototypes learned the hard way:

\begin{itemize}[leftmargin=1.4em]
\item The field named \code{square} stores $q$, not $q^2$. Mixing the two
  conventions silently changes the degree of everything.
\item The weight must be confirmed to be an exact rational \emph{of the right
  type} before it is compared to zero. A weight of $(10^{20}+1)/10^{20}$ offered
  as $1$ is the tested attack.
\item Equality multipliers are unrestricted in sign, so an LP over nonnegative
  variables needs $\pm$ columns for each ideal generator.
\item The target, hypotheses and equations are supplied \emph{from outside} the
  certificate (Invariant~\ref{inv:binding}). The certificate is valid
  \emph{under} those hypotheses; replay does not establish that they hold at any
  particular point, and the Lean adapter must map each hypothesis index to an
  actual proof of the corresponding reified inequality rather than copying a
  polynomial list into the checker.
\end{itemize}

\subsection{Bernstein trees}

\begin{lstlisting}
Tree := { "leaf":  { "degrees": [...], "coefficients": [...] } }
      | { "axis": i, "cut": "1/2", "left": Tree, "right": Tree }
\end{lstlisting}

The root box and the polynomial are separate problem inputs. \emph{Child domains
are computed by the checker during replay}, never read from the certificate, so
a producer cannot omit a strip of the domain by supplying inconsistent
endpoints. Every basis index must appear exactly once in canonical order. A cut
at an interval endpoint, a missing child, or an invalid axis is rejected. Tensor
size is capped \emph{before} expansion, since an adversarial certificate could
otherwise request an enormous one.

A negative Bernstein coefficient is not a negative value of the polynomial; it
is failure of a sufficient condition. Conversely a positive sampled value is not
enough to accept a leaf. Node and depth limits are operational rejection limits,
not mathematical assumptions.

\subsection{Induction traces}

An induction certificate contains the equation, the induction variable, the
\emph{generalised} variables, and four rewrite traces: left and right for the
base case and for the step case. A rewrite step is a
\code{(rule\_identifier, subterm\_path)} pair, optionally with the resulting
term.

The checker derives the base and step contexts and the single permitted
induction hypothesis itself; these are never claims supplied by the certificate.
It reconstructs each substitution and replacement rather than trusting a recorded
output term, and the two replayed sides must end syntactically identical. The
recursion index is a rigid eigenvariable; only the named generalised parameters
may be instantiated. In the base case the induction hypothesis is not available
at all.

The strongest variant attaches to each rule an explicit \emph{flexible} variable
set and requires the trace to supply exactly that substitution domain with
correct sorts.

Certificates are stored and replayed in dependency order, and a lemma is
installed only after its own certificate has been checked. A bundle checker must
additionally reject forward dependencies, shadowing of a defining equation or of
the induction-hypothesis name, and a tampered status field --- which is why an
accepting run takes an explicit list of the entries it is \emph{required} to
have verified, so that an empty or mistyped bundle cannot pass vacuously.

\subsection{Horn derivations}

A node records its ground conclusion, the rule index or an input-fact marker, a
complete ground substitution, and the indices of its premise nodes. Replay
checks groundness; that every referenced rule exists; that the substitution
domain is exactly the head's variables; that the recorded conclusion is that
rule's actual instantiated conclusion; that every premise is already derived;
that every premise index is strictly smaller than the node's own; and that the
last node is the goal.

Trimming retains only ancestors of the root, preserving topological order. A
dependency cycle in the \emph{search} graph may be legitimate; the final
certificate may never justify itself through one.

\subsection{Lattice traces}

Per row: the row number, the unimodular matrix $U$, the gcd $g$, the residual
$r$, and either the quotient $r/g$ or an obstruction marker. Replay recomputes
the current row image and residual, checks that all entries are integers,
verifies $\det U=\pm1$ by independent rational elimination, and performs the
exact state update --- \emph{without} calling the extended-gcd construction or
the solver that produced the trace.

\subsection{SAT and theory proof logs}

Two supported formats.

\emph{RUP form}: a list of added clauses, each re-derived by independent unit
propagation, plus signed negative-cycle theory lemmas validated by exact
coefficient summation with a strictly negative resulting bound. The checker
duplicates the atom-negation rule rather than calling the solver's shared edge
constructor.

\emph{Resolution-DAG form}: nodes tagged \code{input}, \code{theory} or
\code{resolution}, each resolution node carrying parent indices and a pivot.
Cheaper to check, but it replays the solver's own derivation rather than
re-deriving it.

A satisfying assignment is a separate record: it must satisfy every original
clause \emph{and} each atom's truth value must match the integer model.

\subsection{Invariant records}

The record carries the invariant's exact coefficients, the initial state, the
transition, and --- in some prototypes --- the search history and diagnostic
residual strings.

\begin{remark}
Replay must \emph{ignore} the diagnostic residual claims. It independently
reconstructs the invariant, substitutes into the base and step identities, and
compares coefficient maps against zero. This is what prevents a fabricated
\code{"residual": 0} field from serving as evidence.
\end{remark}

\subsection{What the recorded formats actually diverged on}
\label{app:schema-divergence}

For the record, and as the specification for the translation layer:

\begin{longtable}{@{}>{\raggedright\arraybackslash}p{0.26\textwidth}>{\raggedright\arraybackslash}p{0.68\textwidth}@{}}
\toprule
Field & Spellings observed across the nine design-round runs \\
\midrule
\endhead
Polynomial arity &
\code{n} (\src{p1}, \src{p3}, \src{p4}, \src{p6}, \src{p8});
\code{variables} (\src{p2}, \src{p7}); \code{nvars} (\src{p5});
a term-tree AST with no arity field (\src{p9}). \\
Coefficient encoding &
\code{"n/d"} strings in eight runs; \code{\{numerator, denominator\}} integer
objects in \src{p7}. This is the only numeric conversion any adapter needs. \\
Cone container &
\code{terms} (\src{p1}, \src{p3}, \src{p5}); \code{squares} (\src{p2});
\code{nonnegative} (\src{p4}); \code{cone} (\src{p6}); \code{positive}
(\src{p7}, \src{p8}). \\
Weight &
\code{weight} (six runs); \code{coefficient} (\src{p4}, \src{p8}). \\
Hypothesis reference &
\code{powers} (\src{p1}); \code{constraint}, singular and nullable (\src{p2});
\code{hypotheses} (\src{p3}); \code{factors} (\src{p4}, \src{p5}, \src{p6},
\code{p8}); \code{assumptions} (\src{p7}). \\
Ideal multipliers &
\code{equality\_multipliers} (\src{p1}); \code{ideal} (\src{p4}--\src{p7});
\code{equalities} as \code{\{index, multiplier\}} (\src{p8}); absent in
\src{p2}, \src{p3}, \src{p9}. \\
Bernstein node &
\code{\{kind, axis, point, left, right\}} (\src{p2});
\code{\{leaf:\{coefficients, degrees\}\}} (\src{p3});
\code{\{axis, cut, left, right, leaf\}} (\src{p8}). \\
\bottomrule
\end{longtable}

The extension round adds a second, different kind of divergence, and it must not
be resolved the same way.

\begin{longtable}{@{}>{\raggedright\arraybackslash}p{0.26\textwidth}>{\raggedright\arraybackslash}p{0.68\textwidth}@{}}
\toprule
Field & Spellings observed across the extension runs \\
\midrule
\endhead
Closure certificate tag &
\code{pullback-space-v1} and \code{pullback-ideal-v1} (\src{e7});
\code{forge.pullback-space/1} and \code{forge.ideal-space/1} (\src{e8});
\code{check.linear} (\src{e9}); \code{weighted\_closure\_v1} and
\code{inductive\_ideal\_v1} (\src{e6}). \\
Rational encoding &
\code{[numerator, denominator]} pairs nested inside
\code{[exponentVector, rational]} (\src{e7});
flat \code{[exponentVector, numerator, positiveDenominator]} (\src{e8});
bivariate \code{[nExponent, kExponent, numerator, denominator]} (\src{e9}). \\
Negative result &
Separating word (\src{e5}, \src{e6}, \src{e9}); orbit record with an
execution-ordered word (\src{e7}); constructor-tree DAG with an
earlier-child-index discipline (\src{e8}); rooted Kripke model (\src{e4}). \\
Decoder bounds &
Declared per field --- rational bit length $4096$, total degree $128$, $4096$
terms, dimension $64$, at most $32$ transition matrices, order $16$, a
singularity bound cap of $100\,000$, and a $32$\,MiB file cap (\src{e9});
a transition-evaluation cutoff (\src{e8}); a pullback and generator budget
(\src{e1}, \src{e3}). \\
\bottomrule
\end{longtable}

These were \emph{not} unified. A shared tag across \src{e7}, \src{e8} and
\src{e9} would assert that three checkers accept the same object, and they do
not: the three ``space'' certificates differ in whether multipliers are constant
or bounded-degree polynomials, whether the iteration is simultaneous or
round-robin, and whether the initial set is one map's image or a union of
padded images. Renaming keys would hide a mathematical difference rather than a
spelling one. Where the design-round divergence is a translation-layer problem,
this one is a naming problem whose correct resolution is \emph{distinct names}.

One convention from the extension round is worth adopting universally, because
it removes an entire class of forged certificate. \src{e9}'s decoder rejects a
producer-supplied ``the residual is zero'' field outright: no such field is
accepted, and the checker reconstructs the residual from the authoritative
problem and the supplied coefficients, then requires its canonical dictionary to
be empty. The same decoder rejects floats, Booleans masquerading as integers,
non-canonical fractions, non-positive denominators, duplicate monomials and
duplicate JSON keys, with structural shape checks preceding every mathematical
identity. That ordering is the right default for all of the schemas above.

The design-round formats are \emph{semantically compatible and syntactically
incompatible}.
Every one encodes weighted squares times nonnegative-generator products plus
ideal multipliers, over exact rationals. A single canonical schema is reachable
by key renaming plus the one \src{p7} numeric conversion. Three families have no
common schema and stay separate: the CNF/RUP proof logs, the lattice elimination
traces, and the symbolic term-tree theorem bank with its rewrite paths.

\subsection{What a production format must add}

A schema version; a source-goal identifier; an environment and encoding
revision; a variable-map identifier resolved against the Lean side's own
registry; the required assumptions; and the claimed result kind. Resource limits
must bound integer bit sizes, monomial counts, degrees, tree nodes, and proof
depth \emph{before} expensive checking as well as after it.

The Lean process resolves every identifier against its own state. It never
accepts an oracle-provided expression as authoritative, and the decoder itself
must be checked --- a Lean-side decoder with a proved correspondence to the
reified problem is part of the trusted path, not a convenience.
